\usepackage[utf8]{inputenc}
\usepackage[T1]{fontenc}
\usepackage[brazil]{babel}
\usepackage{lmodern}
\usepackage[margin=2.5cm]{geometry}
\usepackage{amsmath,amssymb,amsfonts,bm}
\usepackage{mathtools}
\usepackage{graphicx}
\usepackage{float}
\usepackage{booktabs}
\usepackage{array}
\usepackage{xcolor}
\usepackage{listings}
\usepackage[most]{tcolorbox}
\usepackage{siunitx}
\usepackage[hidelinks]{hyperref}
\usepackage[brazilian,nameinlink]{cleveref}

\sisetup{output-decimal-marker={,}, group-separator={.}}

% ---------- Layout: tabelas compactas e quebra de nomes longos ----------
\emergencystretch=3em
\setlength{\tabcolsep}{4pt}
\AtBeginEnvironment{tabular}{\small}
\usepackage{xurl}
% Ordinais (a fonte Latin Modern no XeTeX/Tectonic não tem "ª"/"º")
\usepackage{newunicodechar}
\newunicodechar{ª}{\textsuperscript{a}}
\newunicodechar{º}{\textsuperscript{o}}
\graphicspath{{figuras/}}

% ---------- Caixas ----------
\newtcolorbox{dica}[1][]{colback=blue!4,colframe=blue!50!black,
  fonttitle=\bfseries,title={Dica prática},#1}
\newtcolorbox{receita}[1][]{colback=green!4,colframe=green!40!black,
  fonttitle=\bfseries,title={Receita de projeto},#1}
\newtcolorbox{codigo}[1][]{colback=orange!4,colframe=orange!60!black,
  fonttitle=\bfseries,title={No código (C++)},#1}
\newtcolorbox{atencao}[1][]{colback=red!4,colframe=red!60!black,
  fonttitle=\bfseries,title={Atenção},#1}

% ---------- Listagens MATLAB ----------
\definecolor{mlkw}{RGB}{0,0,200}
\definecolor{mlcm}{RGB}{34,139,34}
\definecolor{mlst}{RGB}{160,32,240}
\lstdefinestyle{matlab}{language=Matlab,basicstyle=\ttfamily\footnotesize,
  keywordstyle=\color{mlkw},commentstyle=\color{mlcm},stringstyle=\color{mlst},
  numbers=left,numberstyle=\tiny\color{gray},frame=single,breaklines=true,
  showstringspaces=false,columns=fullflexible,
  literate={á}{{\'a}}1 {ã}{{\~a}}1 {é}{{\'e}}1 {ê}{{\^e}}1 {í}{{\'i}}1
           {ó}{{\'o}}1 {õ}{{\~o}}1 {ú}{{\'u}}1 {ç}{{\c c}}1 {â}{{\^a}}1}
\lstset{style=matlab}

% ---------- Notação ----------
\newcommand{\R}{\mathbb{R}}
\newcommand{\vect}[1]{\bm{#1}}
\newcommand{\ep}{\vect{e}_p}
\newcommand{\ev}{\vect{e}_v}
\newcommand{\rr}{\vect{r}}
\newcommand{\rhat}{\hat{\vect{r}}}
\newcommand{\T}{^{\mathsf{T}}}
\newcommand{\Ts}{T_s}
\newcommand{\code}[1]{\texttt{#1}}
% Permite quebrar linha depois de "_" (identificadores longos do C++/ROS)
\let\underscoreorig\_
\renewcommand{\_}{\underscoreorig\allowbreak}
\DeclareMathOperator{\sat}{sat}
\DeclareMathOperator{\sgn}{sgn}
